% Problem Set 8 -- Part 1: question (f). Numbers from ps8_partA.py (part_f), using the monthly
% returns of the frozen strategy from question_e/part_e.ipynb.

\subsection*{Question (f): confidence intervals for the Sharpe ratio}

We use the monthly excess returns of the frozen strategy from (e), both in the walk-forward
training period (1953:03--2001:06, $T=580$) and in the held-out test period (2001:07--2026:06,
$T=300$). Let $\widehat{SR}$ be the monthly Sharpe ratio, $\hat\gamma_3$ the skewness and
$\hat\gamma_4$ the (raw, not excess) kurtosis of the returns. Under i.i.d.\ normal returns
(Lo, 2002) the Sharpe ratio is approximately normal with
\[
  \operatorname{Var}\big(\widehat{SR}\big)\approx\frac{1+\tfrac12\widehat{SR}^{\,2}}{T},
\]
and for i.i.d.\ returns with any distribution (Mertens, 2002; Opdyke, 2007) the variance becomes
\[
  \operatorname{Var}\big(\widehat{SR}\big)\approx\frac{1-\hat\gamma_3\,\widehat{SR}
  +\tfrac{\hat\gamma_4-1}{4}\,\widehat{SR}^{\,2}}{T}.
\]
The 95\% interval is $\widehat{SR}\pm1.96\,\mathrm{SE}$, computed in monthly units and then
multiplied by $\sqrt{12}$. Table~\ref{tab:srci} reports both intervals.

\begin{table}[H]
\centering\small
\caption{Annualized Sharpe ratios of the strategy from (e) with 95\% confidence intervals. ``Normal'' assumes i.i.d.\ normal returns; ``General'' adjusts for skewness and kurtosis. Net returns subtract the trading and borrowing costs used in (e).}
\label{tab:srci}
\begin{tabular}{llccccc}
\toprule
Sample & Returns & Skew & Kurtosis & $\widehat{SR}$ (ann.) & Normal 95\% CI & General 95\% CI \\
\midrule
Training, 1953--2001 & Gross & $-$0.03 & 11.0 & 1.53 & [1.24, 1.83] & [1.19, 1.88] \\
($T=580$)            & Net   & $-$0.05 & 11.1 & 1.34 & [1.04, 1.63] & [1.00, 1.67] \\
\addlinespace
Test, 2001--2026     & Gross & 0.02 & 5.8 & 0.34 & [$-$0.05, 0.74] & [$-$0.05, 0.74] \\
($T=300$)            & Net   & 0.04 & 5.8 & 0.11 & [$-$0.28, 0.50] & [$-$0.28, 0.50] \\
\bottomrule
\end{tabular}
\end{table}

\begin{tight}
  \item \textbf{The training Sharpe ratio is clearly positive under both distributions.} The net
  training Sharpe ratio of 1.34 has a 95\% interval of roughly 1.0 to 1.7, so even its lower end
  is far above zero and above the market's long-run Sharpe ratio of about 0.4--0.5. Taken at face
  value, this would say the strategy is a very strong one.

  \item \textbf{The general interval is wider, and the gap grows with the Sharpe ratio.} The
  strategy's returns are nearly symmetric but have fat tails, with kurtosis of about 11 in
  training against 3 for a normal distribution. Fat tails make the estimated Sharpe ratio
  noisier, so the general interval is about 15\% wider than the normal one in training. In the
  test period the adjustment barely matters, because the kurtosis term is multiplied by
  $\widehat{SR}^{\,2}$, which is close to zero there.

  \item \textbf{The test interval includes zero, so the training interval is too optimistic.} The
  held-out net Sharpe ratio of 0.11 falls far below the lower end of the training interval and
  below the market's 0.58 over the same years, and its own interval runs from about $-0.3$ to
  $0.5$. Both intervals assume the strategy was fixed before we saw the data, but the signals in
  (e) were chosen using facts from (a), (b) and (d) over the full sample, and the rule was
  revised once after earlier test results were seen. Autocorrelation in the returns is small (about 0.07 at
  one lag), so it does not explain the gap. The training interval therefore measures sampling
  noise only, not the selection bias that part (g) is about, which
  is why the strategy looks much weaker once it is run on data it has not seen.
\end{tight}
